% 03_metabolic.tex
\begin{landscape}

	% =============================================================================
	% CONFIGURACIÓN
	% =============================================================================
	\setlength{\tabcolsep}{2pt}
	\setlength{\extrarowheight}{3pt}
	\fontsize{10pt}{10pt}\selectfont

	% =============================================================================
	% CÁLCULO DE ANCHOS DE COLUMNAS
	% =============================================================================
	\setlength{\availablewidth}{\linewidth}

	% Restamos (8 columnas * 2 lados) = 16 espacios para un ajuste matemático exacto
    \addtolength{\availablewidth}{-16\tabcolsep}

    % Restamos el ancho de las líneas verticales (9 líneas * 0.4pt estándar)
    \addtolength{\availablewidth}{-3.6pt}

	\newcommand{\tableheader}{\rowcolor{blue!10}
		\textbf{ESPE\-CIE} & \textbf{Nº DE CASOS} & \textbf{ENFERME\-DA\-DES} & \textbf{FÁRMACO} &
		\textbf{PRINCIPIO ACTIVO} & \textbf{DOSIS/VÍA DE ADMINISTRACIÓN} &
		\textbf{DURA\-CIÓN} & \textbf{OBSERVA\-CIO\-NES} \\
	}

	% =============================================================================
	% LONGTABLE CON CAPTION AL PRINCIPIO
	% =============================================================================
	\begin{longtable}{
		|C{0.07\availablewidth}           % ESPECIE
		|C{0.07\availablewidth}           % Nº CASOS
		|C{0.11\availablewidth}           % ENFERMEDADES
		|L{0.12\availablewidth}          % FÁRMACO
		|J{0.3\availablewidth}          % PRINCIPIO ACTIVO
		|L{0.17\availablewidth}          % DOSIS/VÍA
		|C{0.06\availablewidth}            % TIEMPO
		|L{0.1\availablewidth}|         % OBSERVACIONES
	}

		% ──────────────── CAPTION COMO TÍTULO (solo en primera página) ────────────────
		\caption{\raggedright Registro de Tratamientos para Enfermedades Metabólicas}
		\label{tab:metabolicas} \\
		\hline
		\tableheader
		\hline
		\endfirsthead

		% ──────────────── ENCABEZADO PARA PÁGINAS SIGUIENTES ────────────────
		\multicolumn{8}{l}{\normalsize\textbf{(Continuación) Cuadro \ref{tab:metabolicas}}. Registro de Tratamientos para Enfermedades Metabólicas} \\[6pt]
		\hline
		\tableheader
		\hline
		\endhead

		% ──────────────── PIE DE PÁGINA (continuación) ────────────────
		\hline
		\multicolumn{8}{r}{\footnotesize\textit{Continúa en la siguiente página \ldots}}
		\endfoot

		% ──────────────── PIE FINAL ────────────────
		%\hline
		\multicolumn{8}{l}{\footnotesize\textit{\textbf{Fuente}: Elaboración propia.}}
		\endlastfoot

		% =============================================================================
		% CONTENIDO DE LA TABLA (extraído del PDF '03_metabolica.pdf')
		% =============================================================================

		% Canino – Ascitis
		\multirow{3}{*}{\textbf{Canino}} & \multirow{3}{*}{3} & \multirow{3}{*}{Ascitis} &
		Furozur & Furosemida, dexametasona, fosfato de sodio & 1ml/10kg IM & \multirow{3}{*}{3-5 días} & \multirow{3}{*}{Satisfactorio} \\
		\cline{4-6}
		& & & Cefotaxima & Cefotaxima sódica & 2,5ml/10kg IV-IM & & \\
		\cline{4-6}
		& & & AminoLab forte & Aminoácidos, dextrosa, cloruro de potasio, cloruro de sodio, cloruro de calcio, acetato de sodio, vitamina B1 B2 B6 B12 & 1ml/1kg IV & & \\
		\hline

		% Canino – Anemia
		\multirow{4}{*}{\textbf{Canino}} & \multirow{4}{*}{2} & \multirow{4}{*}{Anemia} &
		Ringer Lactato & Electrolitos & Según criterio médico IV & \multirow{4}{*}{3-5 días} & \multirow{4}{*}{Satisfactorio} \\
		\cline{4-6}
		& & & AminoLab forte & Aminoácidos, dextrosa, cloruro de potasio, cloruro de sodio, cloruro de calcio, acetato de sodio, vitamina B1 B2 B6 B12 & 1ml/1kg IV & & \\
		\cline{4-6}
		& & & Complejo B & Vitamina B1, Vitamina B2, Vitamina B6, Vitamina B12, Vitamina B15 & 1ml/20kg IM-SC & & \\
		\cline{4-6}
		& & & Duplafer & Complejo hierro dextano, vitamina B12 & 1-3ml IM & & \\
		\hline

		% Felino – Anemia
		\multirow{4}{*}{\textbf{Felino}} & \multirow{4}{*}{1} & \multirow{4}{*}{Anemia} &
		Ringer Lactato & Electrolitos & Según criterio médico IV & \multirow{4}{*}{3-5 días} & \multirow{4}{*}{Satisfactorio} \\
		\cline{4-6}
		& & & AminoLab forte & Aminoácidos, dextrosa, cloruro de potasio, cloruro de sodio, cloruro de calcio, acetato de sodio, vitamina B1 B2 B6 B12 & 1ml/1kg IV & & \\
		\cline{4-6}
		& & & Complejo B & Vitamina B1, Vitamina B2, Vitamina B6, Vitamina B12, Vitamina B15 & 1ml/20kg IM-SC & & \\
		\cline{4-6}
		& & & Duplafer & Complejo hierro dextano, vitamina B12 & 1-3ml IM & & \\
		\hline

		\rowcolor{green!10} % Verde claro para el total
		\textbf{TOTAL} & \textbf{6} \\
		\cline{1-2}

	\end{longtable}
	\vspace{-12pt}\normalsize
	% Resumen conciso de la tabla
	El cuadro registra los tratamientos aplicados para enfermedades metabólicas en caninos y felinos durante el internado. Se documentaron 6 casos (3 ascitis, 3 anemia). Los tratamientos incluyen diuréticos, antibióticos, fluidoterapia, aminoácidos, vitaminas y suplementos de hierro, administrados IV/IM/SC durante 3-5 días, con resultados satisfactorios en todos los casos.

\end{landscape}
